\chapter{Vitiligo}

\section*{Definition}
loss of skin pigment is a chronic autoimmune disorder characterized by depigmented, well-circumscribed macules and patches resulting from selective destruction of epidermal melanocytes. It affects 0.5--2\% of the global population regardless of sex or race, with onset before age 20 in half of cases.

\section*{Pathophysiology}
CD8+ cytotoxic T lymphocytes infiltrate the epidermis and specifically target melanocytes through recognition of melanocyte-specific antigens (MART-1, gp100, tyrosinase). Interferon-\ensuremath{\gamma} and IL-15 drive T-cell activation and persistence. Keratinocytes in loss of skin pigment skin overexpress CXCL10, maintaining T-cell recruitment. Oxidative stress with accumulation of reactive oxygen species and defective antioxidant responses precedes melanocyte destruction.

\section*{Causes}
\begin{itemize}
  \item \textbf{Autoimmune:} Associated with autoimmune thyroid disease (Hashimoto's, Graves'), pernicious anemia, Addison's disease, type 1 diabetes, hair loss areata
  \item \textbf{Genetic:} Polymorphisms in HLA, PTPN22, NLRP1, FOXP3, and TYR genes; family history in 15--20\%
  \item \textbf{Triggers:} Physical trauma (Koebner phenomenon), sunburn, emotional stress, chemical exposure (phenols, catechols)
  \item \textbf{Subtypes:} Segmental (unilateral, dermatomal, early onset, stable after 1--2 years) and non-segmental (bilateral, symmetrical, progressive)
\end{itemize}

\section*{When to See a Doctor}
See your doctor if:
\begin{itemize}
  \item New or progressive depigmented patches, especially in cosmetically sensitive areas (face, hands, genitalia)
  \item Premature graying of scalp or body hair (loss of skin pigment-associated poliosis)
  \item Evaluate for associated autoimmune conditions: fatigue, weight changes, palpitations, cold intolerance
  \item Consider referral for phototherapy or systemic treatment if extensive or rapidly progressive
\end{itemize}

\section*{Related Symptoms}
\begin{itemize}
  \item Leukoderma, piebaldism, tinea versicolor, pityriasis alba, post-inflammatory hypopigmentation
\end{itemize}
\textbf{Important:} Segmenter vitiligo is more common in children and typically stabilizes within 1--2 years, whereas non-segmental vitiligo tends to have a relapsing--remitting or slowly progressive course throughout life.

\section*{Self-Care / Home Management}
\begin{itemize}
  \item Broad-spectrum sunscreen (SPF 50+) to depigmented areas to prevent sunburn, which also triggers Koebner phenomenon
  \item Cosmetic camouflage with waterproof, transfer-resistant makeup or self-tanning agents (dihydroxyacetone)
  \item Wear protective clothing and avoid known physical triggers (friction, trauma)
\end{itemize}

\section*{How Your Doctor Will Evaluate You}
\begin{itemize}
  \item Clinical diagnosis: Well-demarcated, non-scaly, depigmented macules; Wood's lamp examination enhances contrast
  \item Skin biopsy shows complete absence of melanocytes and melanin in affected skin; lymphocytic infiltrate at active border
  \item Screen for autoimmune thyroid disease: TSH, free T4, anti-TPO antibodies at diagnosis and annually
\end{itemize}

\section*{Treatment Options}
\begin{itemize}
  \item First-line: Topical corticosteroids (mometasone, betamethasone) or calcineurin inhibitors (tacrolimus, pimecrolimus) for limited disease
  \item Narrowband UVB phototherapy 2--3 times weekly for moderate--severe or rapidly progressive disease
  \item Excimer laser (308 nm) for localized, recalcitrant patches
  \item Systemic: Oral corticosteroids (mini-pulse), methotrexate, or JAK inhibitors (ruxolitinib cream, oral tofacitinib) for aggressive disease
  \item Surgical: Autologous melanocyte transplantation for stable, segmental, or treatment-resistant disease
\end{itemize}

\clearpage
